\section{Pair Mutual Correlated Agreement}

\subsection{Bad challenges}
\label{sec:mca-bad-challenges}

Let \(p=(p_0,p_1)\in(\F^2)^I\), and write
\[
  y_\gamma=p_0+\gamma p_1.
\]
For \(S\subseteq I\), say that \(S\) is \(\delta\)-large when
\[
  \frac{|S|}{n}\ge1-\delta.
\]
Define \(\Bad_{\MCA}(\mathcal C,\delta,p)\) to be the set of
\(\gamma\in\F\) for which there exists a \(\delta\)-large set \(S\) such that:
\begin{enumerate}[label=(\roman*)]
  \item \(y_\gamma\) agrees on \(S\) with a codeword of \(\mathcal C\);
  \item the pair \(p\) does not agree on the same \(S\) with a codeword of
    \(\Cpair\).
\end{enumerate}
For a fixed pair \(p\) and a uniform draw \(\gamma\in\F\), the failure
probability is \(|\Bad_{\MCA}(\mathcal C,\delta,p)|/|\F|\).
Maximizing over \(p\) gives the MCA error.
\[
  \varepsilon_{\MCA}(\mathcal C,\delta)
  =
  \max_{p\in(\F^2)^I}
  \frac{|\Bad_{\MCA}(\mathcal C,\delta,p)|}{|\F|}.
\]

This is the two-word same-set event appearing in
\cite[Theorem~4.6, full version]{benSassonCarmonHabockKoppartySaraf2026}.
That source uses polynomials of degree at most its parameter, so its
parameter is \(k-1\) in the degree-\(<k\) convention used here.
Our bound below is uniform in the radius; the cited literature gives
sharper radius-dependent bounds in its stated regimes. At fixed rate and a
fixed positive gap below the Johnson radius, the cited theorem gives an
exceptional-challenge count linear in \(n\), generally much smaller than
the binomial bound below. Broader polynomial-generator formulations of MCA
for linear codes are studied by Bordage et al.
\cite{bordageEtAl2026polynomialMCA}.
Jo's local-test count for MDS codes implies the uniform binomial bound
below and also gives a sharper radius-dependent bound
\cite[Theorem~4.2]{jo2026beyondJohnson}. We give the RS interpolation proof
underlying our Lean certificate map.

\begin{theorem}[Uniform pair-MCA certificate bound]
\label{thm:pair-mca}
Let \(\mathcal C=\RS[\F,\mathcal L,k]\) have injectively indexed
evaluation domain of length \(n\), where \(1\le k\le n\).  For every
\(\delta\in\mathbb R\) and every \(p\in(\F^2)^n\),
\[
  |\Bad_{\MCA}(\mathcal C,\delta,p)|
  \le \binom n{k+1}.
\]
Consequently,
\[
  \varepsilon_{\MCA}(\mathcal C,\delta)
  \le \frac{\binom n{k+1}}{|\F|}.
\]
Smoothness is not required for this bound.
\end{theorem}

\begin{proof}
Fix a bad challenge \(\gamma\), and choose a witnessing set \(S\).
Necessarily \(|S|>k\).  Indeed, when \(|S|\le k\), interpolate each component
of \(p\) on \(S\) by a polynomial of degree \(<k\); the resulting pair
codeword agrees with \(p\) on \(S\), contradicting condition~(ii).

Choose a \(k\)-subset \(T\subset S\).  Let \(a_T,b_T\) be the unique
degree-\(<k\) polynomials interpolating \(p_0,p_1\), respectively, on \(T\).
Condition~(ii) says that the pair interpolant \((a_T,b_T)\) fails somewhere
on \(S\). Choose such an \(i\in S\setminus T\), and set
\(U=T\cup\{i\}\). Condition~(i) remains true after restriction from \(S\) to
\(U\). The pair still cannot agree with \(\Cpair\) on \(U\): any purported
pair codeword would agree with \((a_T,b_T)\) on \(T\), hence equal it by
\(k\)-point determination, contradicting the mismatch at \(i\).
Thus \(U\) is a \((k+1)\)-coordinate obstruction for \(\gamma\).

A fixed \(U\) obstructs at most one challenge. Suppose distinct
\(\gamma,\eta\) were both bad on \(U\), with scalar codewords
\(c_\gamma,c_\eta\) agreeing with
\(p_0+\gamma p_1,p_0+\eta p_1\), respectively. By linearity,
\[
  d=\frac{c_\gamma-c_\eta}{\gamma-\eta},
  \qquad
  a=c_\gamma-\gamma d
\]
are codewords. On \(U\), they agree with \(p_1\) and \(p_0\), respectively,
so \((a,d)\) is a pair codeword agreeing with \(p\) on \(U\), a
contradiction.

Make the finite choices deterministic. Mapping each bad challenge to its
unordered obstruction \(U\) is injective, and there are
\(\binom n{k+1}\) such subsets.
Dividing by \(|\F|\) proves the error bound.  Notice that the large-set
inequality is no longer used after the witness is extracted; this is why the
bound is uniform in \(\delta\).
\end{proof}

\begin{corollary}[Explicit smooth family at \(2^{-128}\)]
\label{cor:explicit-family-mca}
For every member \(\mathcal C_u\) of the family in
\cref{sec:explicit-family}, every \(\delta\in[0,1]\), and every pair word
\(p\),
\[
  \frac{|\Bad_{\MCA}(\mathcal C_u,\delta,p)|}{|\F_u|}
  \le 2^{-128}.
\]
Thus, for this family and this same-set pair-MCA definition, every normalized
radius is safe and the largest admissible radius is \(1\).
\end{corollary}

\begin{proof}
\Cref{sec:field-inequality} proves
\[
  2^{128}\binom{n_u}{k_u+1}\le q_u.
\]
Apply \cref{thm:pair-mca}.
\end{proof}

\subsection{An exact small-radius boundary instance}
\label{sec:mca-boundary-field}

Jo \cite{jo2026beyondJohnson} gives exact 128-bit MCA thresholds for four
length-\(64\) smooth codes. Chojecki \cite{chojecki2026shortening,chojecki2026thresholds}
gives exact finite small-radius counts and beyond-Johnson certificates.
The example below certifies the field arithmetic, domain, and adjacent
threshold witnesses for a length-\(4096\) code in Lean.

The uniform certificate bound is sufficient but need not locate the first
unsafe radius.  Write
\[
  B_{\mathcal C}(r)
  =
  \max_p|\Bad_{\MCA}(\mathcal C,r/n,p)|
\]
for the exact worst bad-challenge count at count radius \(r\).  Consider
\[
  q=2^{128}+192513
   =340282366920938463463374607431768403969.
\]
The accompanying Lean development certifies that \(q\) is prime and
\(q\equiv1\pmod {4096}\).  Hence \(\mathbb F_q^\times\) contains an
order-\(4096\) subgroup, which we use as the evaluation domain for the smooth
RS code with \(n=4096\) and \(k=1024\).

\begin{theorem}[Exact 128-bit small-radius transition]
\label{thm:mca-boundary-threshold}
For this code,
\[
  B_{\mathcal C}(0)=1,
  \qquad
  B_{\mathcal C}(1)\ge2.
\]
Consequently, for the \(2^{-128}\) MCA error target, its greatest safe count
radius is \(0\) and its first unsafe count radius is \(1\).
For the relative-radius definition over
all \(\delta\in[0,1]\), the safe set is
\[
  \left[0,\frac1{4096}\right),
\]
so the real-radius supremum is \(1/4096\) and is not attained.
\end{theorem}

\begin{proof}
At count radius zero, every large agreement set is the whole coordinate set.
A fixed set supports at most one bad challenge, while a unit-spike affine
line realizes one; hence \(B_{\mathcal C}(0)=1\).  At count radius one, take
two coordinates outside a \(k\)-point interpolation set.  The affine line
from the unit spike at the first coordinate to the unit spike at the second
has distinct bad challenges \(0\) and \(1\), proving
\(B_{\mathcal C}(1)\ge2\).

The certified field lies in the exact budget window
\[
  2^{128}\le q<2^{128}\cdot2.
\]
Thus the integer criterion \(2^{128}B_{\mathcal C}(r)\le q\) holds at
\(r=0\) and fails at \(r=1\).  Monotonicity of the bad-count spectrum gives
the two grid endpoints.

Finally, if \(0\le\delta<1/n\), the inequality
\(|S|/n\ge1-\delta>1-1/n\) forces the integer \(|S|\) to equal \(n\).
The bad sets are therefore exactly those at count radius zero throughout
that half-open interval.  At \(\delta=1/n\), sets missing one coordinate
become admissible, and the two-spike line already violates the budget.
Monotonicity preserves failure for larger radii, proving the real-radius
claim.
\end{proof}

\subsection{Distance interpretation}

Let
\[
  D_2(p)=\Delta(p,\Cpair),
  \qquad
  D_1(y)=\Delta(y,\mathcal C).
\]
For a challenge outside the MCA bad set, scalar closeness on a large agreement
set forces pair closeness on that same set. For \(\delta\in[0,1]\), this
gives the distance implication
\[
  D_2(p)\ge\delta
  \Longrightarrow
  D_1(y_\gamma)\ge\delta
\]
for every MCA-good \(\gamma\), with the fraction of excluded challenges bounded
by \(\varepsilon_{\MCA}(\mathcal C,\delta)\).

To see the boundary convention explicitly, suppose that
\(D_1(y_\gamma)<\delta\), and choose a closest scalar codeword \(c\).
Let \(S=\{i:y_\gamma(i)=c(i)\}\). Then
\(|S|/n=1-D_1(y_\gamma)>1-\delta\), so \(S\) is \(\delta\)-large.
Since \(\gamma\) is MCA-good, a pair codeword agrees with \(p\) on
\(S\). Hence \(D_2(p)\le D_1(y_\gamma)<\delta\). Taking the
contrapositive proves the displayed implication, including its non-strict
\(\ge\) boundary.

This is a challenge-indexed statement.  Replacing the sampled challenge by one
fixed field element would require a separate theorem that the selected element
is good for every pair, which does not follow from a fractional MCA bound.

\subsection{Endpoint interpretation}

At \(\deltacap=1-k/n\), interpolation on any \(k\) coordinates shows that
\[
  D_1(y)\le\deltacap
  \quad\text{and}\quad
  D_2(p)\le\deltacap
\]
for every scalar word \(y\) and pair word \(p\).  Moreover, on every set
\(S\) of size at most \(k\), each component of \(p\) extends to a
degree-\(<k\) polynomial.  Such an \(S\) can therefore never witness an MCA
failure.  Every bad endpoint challenge must instead have a witness with
\(|S|\ge k+1\).

Consequently, the endpoint distance implication above is an equality
non-collapse statement: when \(D_2(p)=\deltacap\), a good challenge must not
lower the projected distance strictly below \(\deltacap\).
\Cref{thm:pair-mca} supplies a uniform quantitative bound on the
excluded challenges.

\subsection{Commit-then-challenge order}

The experiment fixes both components of \(p\) before sampling
\(\gamma\) uniformly from \(\F\). The witnessing set \(S\) may depend on
the draw. A protocol using this bound must bind the entire pair before
deriving its challenge.

\subsection{Scope}

For a supplied field, the certificate proof gives the sufficient comparison
\[
  \binom n{k+1}\le 2^{-128}|\F|.
\]
Relative to the older ordered-certificate count used in this development,
the bound saves the exact factor \(k+1\), since
\(\binom nk(n-k)=(k+1)\binom n{k+1}\).
This is not a uniform improvement over the cited literature bounds.
If that inequality fails, this proof does not determine the optimal radius or
show that the MCA target is impossible; it only ceases to certify the requested
error.

For an arbitrary supplied finite code, radius \(r/n\) meets the
\(2^{-128}\) target exactly when
\[
  2^{128}B_{\mathcal C}(r)\le|\F|.
\]
This is also the exact formal equivalence: the Lean theorem identifies the
real-valued operational error predicate at normalized radius \(r/n\) with the
displayed integer inequality, so no rounding relaxation is hidden here.
Thus a uniform upper-bound certificate at \(r\), together with one concrete
affine line exceeding the budget at \(r+1\), pins both the largest safe and
first unsafe Hamming-grid radii. The remaining middle regime is the
code-specific problem of proving those matching bounds; smoothness by itself
does not determine this incidence spectrum.

The result is for the two-coordinate, same-set MCA experiment defined above.
Jo \cite{jo2026interleaving} proves that the affine-line MCA error itself
is invariant under row-wise interleaving.

The smooth extension-field family \(q_u=3^{n_u}\) meets the list
budget at capacity and the MCA budget at every normalized radius. In the
prime field of size \(2^{128}+192513\), the exact MCA transition is between
count radii zero and one, and the list transition is between \(1536\) and
\(1537\) (\cref{cor:boundary-field-list}). In the canonical application field of size
\(12289\), the radius-zero unit-spike line already has one bad challenge.
Since \(12289<2^{128}\), this last code has no safe count radius on its
full pair space.
\Cref{sec:jo-paired-thresholds,sec:subfield-crossings} gives additional
exact crossings on the length-\(64\) and length-\(128\) domains.

\subsection{Projective incidence of bad challenges}
\label{sec:mca-projective-incidence}

Write \(q=|\F|\). Let \(V=\F^I/\mathcal C\) be the syndrome space and let
\(H:\F^I\to V\) be the quotient map. For an error support \(E\subseteq I\), put
\(U_E=\operatorname{span}\{He_i:i\in E\}\), where \(e_i\) is the coordinate
vector at site \(i\). A word agrees with a codeword on \(I\setminus E\)
exactly when its syndrome belongs to \(U_E\). Syndrome-space formulations
also appear in the study of correlated agreement for random codes
\cite{yuanZhu2026syndromeSpace}. The event in \cref{sec:mca-bad-challenges}
uses the same support \(E\) in the scalar and pair tests.

Write \(s_j=Hp_j\) for the two syndromes of a received pair. At count
radius \(r\), a challenge \(\gamma\) is bad exactly when some
\(|E|\le r\) satisfies
\[
  s_0+\gamma s_1\in U_E,\qquad
  \operatorname{span}(s_0,s_1)\nsubseteq U_E.
\]
The second condition excludes the case where both columns admit codeword
comparisons on \(I\setminus E\). This is the full-field case of Chojecki's
syndrome-line criterion \cite[Section~3]{chojecki2026shortening}.
If the two syndromes span a space of
dimension at most one, there is at most one bad challenge. Indeed, any
nonzero combination generates their common line in the dimension-one
case, so only a combination with zero syndrome can witness failure. A
pair of zero syndromes satisfies both component agreement tests everywhere.

If the syndromes are independent, the plane \(W=\operatorname{span}(s_0,s_1)\)
defines a projective line \(\mathbb P(W)\). Each \(E\) with
\(\dim(W\cap U_E)=1\) contributes one point. Count coincident points once
and exclude the point \([s_1]\) at infinity. The affine coordinates
\([s_0+\gamma s_1]\) of the remaining points give exactly the bad
challenges. Write \(N_r(W)\) for the number of distinct contributed points
before excluding infinity. Every basis of \(W\) lifts to a received pair,
and its second vector can be chosen on any rational point of
\(\mathbb P(W)\). Thus the maximum affine count over bases of \(W\) is
\(\min(q,N_r(W))\).

\paragraph{Equations for multiple hits.}
For \(W\in\operatorname{Gr}(2,V)\), let \(\xi\) be a nonzero Pl\"ucker
representative and let \(\pi_E:V\to V/U_E\). The incidence
\(W\cap U_E\ne0\) is the linear equation in Pl\"ucker coordinates
\[
  (\wedge^2\pi_E)(\xi)=0.
\]
Remove the closed locus of planes contained in \(U_E\) to keep a single
intersection point. For supports \(E,F\) satisfying this rank condition,
their points are distinct exactly when
\[
  \bigl(\wedge^2(V\to V/(U_E\cap U_F))\bigr)(\xi)\ne0.
\]
Indeed the points coincide exactly when their common line lies in both
syndrome subspaces. These tests use the intersection of the syndrome
subspaces themselves.

For a bad-count budget \(1\le T<q\), an upper bound of \(T\) is equivalent
to the absence of a rational plane with \(T+1\) distinct contributed
points. This gives a finite Gr\"obner-basis test on the Pl\"ucker charts.
For each choice of supporting sets, impose the Grassmannian equations
and the linear incidence equations, and encode the open conditions by
inverse witnesses on charts where a specified coordinate is nonzero.
Adjoining \(X^q-X=0\) for every variable makes the unit-ideal test a
certificate of emptiness over \(\F\). Section~8 shows how an elimination
closure can add points with empty fibers.

\paragraph{A four-point example.}
For \(\mathcal C=\RS[\mathbb F_5,\mathbb F_5^\times,2]\) at \(r=1\),
take a parity-check matrix with column \((a,a^2)^{\mathsf T}\) at each
\(a\in\mathbb F_5^\times\). The sums of the first three positive powers
over the domain vanish, so the matrix annihilates affine evaluation
words. Its rank is two. The singleton supports give four distinct
directions \([1:a]\). Choosing \(s_0=(1,0)\), \(s_1=(0,1)\) places infinity
outside those four directions, giving bad challenges \(1,2,3,4\). The
four directions bound the count for independent syndromes, and the
dependent case contributes at most one. Hence \(B_{\mathcal C}(1)=4\).

\paragraph{An empty matching over \(\mathbb F_{17}\).}
For the smooth \([8,4]\) code at the powers of \(9\) in \(\mathbb F_{17}\),
the parity-check column at \(a\) is \((a,a^2,a^3,a^4)^{\mathsf T}\).
For each of the supports \(\{1,9\},\{13,15\},\{16,8\},\{4,2\}\), the
coefficients of \((X-a)(X-b)\) and \(X(X-a)(X-b)\) give two annihilating rows.
The linear incidence equations in order \(01,02,03,12,13,23\) give
\((13v,0,5u,u,0,v)\). The Pl\"ucker relation \(5u^2+13v^2=0\) gives
\(v\ne0\) for any projective solution and \(x^2=11\) for \(x=u/v\).
Set \(S(X)=\sum_{j=0}^{7}11^j X^{14-2j}\). In \(\mathbb F_{17}[X]\),
\[
  1=(7X^2S(X)+3)(X^2-11)+10X(X^{17}-X).
\]
Lean checks the reduction and this identity. Together with the field
equation \(X^{17}-X=0\), they exclude the specified matching.

\paragraph{The Frobenius orbit of the matching.}
The coincidence locus of two support labels is the pullback of the
diagonal under their contributed-point maps. Deleting those loci keeps
only distinct points in the marked incidence space. If a resulting
locally closed locus is finite and reduced over \(K=\mathbb F_b\), its
geometric points form a finite set with a Frobenius action
\cite{stacksEtalePoint}. An orbit of size \(d\) contributes \(d\) rational
points over \(\mathbb F_{b^e}\) exactly when \(d\mid e\). For a pencil
already rational over \(K\), every one-dimensional support intersection
is rational over \(K\) by linear algebra. Extension changes which pencils
can occur.

For the four supports above, adjoining \(x\) with \(x^2=11\) gives
the independent syndrome columns
\[
 s=(1,0,13x,0),\qquad t=(0,13,0,5x).
\]
Their Pl\"ucker vector is \((13,0,5x,x,0,1)\). The four affine
parameters in the listed support order are \(7+12x,6+14x,10+5x,11+3x\).
The vector \(s+\gamma t\) lies in the corresponding support span at
each parameter, and \(t\) lies in none of those four spans. The parameters
are distinct. The two choices of \(x\) give all geometric solutions by
the reduced chart, and \(x^{17}=x\cdot 11^8=-x\). Thus the two pencils
form one Frobenius orbit, and the matching is realizable over
\(\mathbb F_{17^e}\) exactly for even \(e\).
Its normalized chart has coordinate algebra
\(\mathbb F_{17}[X]/(X^2-11)\), a finite separable quadratic field.

\paragraph{The finite \texorpdfstring{\'etale}{etale} part of the matching census.}
Set \(K=\mathbb F_{17}\) and \(K_2=\mathbb F_{17^2}\). A perfect matching
partitions the eight labelled sites into four unordered pairs. There are
\(7\cdot5\cdot3=105\) matchings, by selecting the partner of the least
remaining site at each step. For each matching \(M\), stack the four
signed-minor rows of the annihilator pairs to give a matrix \(L_M\) on
the six Pl\"ucker coordinates. With \(K_M=\ker L_M\), its incidence locus
is the zero scheme of the Grassmann quadratic restricted to
\(\mathbb P(K_M)\).

Any four parity columns are independent. Thus different pairs in a
matching have disjoint vector spans, and a plane meeting all four spans
has a different one-dimensional intersection with each. This holds over
the algebraic closure, so the open conditions for distinct single hits
hold on the whole matching locus. Over a finite extension, the four hits
are rational on each rational plane, and the affine chart can place its
point at infinity outside them.

Row reduction over \(K\) gives the following census. For a two-dimensional
kernel, the discriminant of the restricted binary quadratic gives the
first three cases. For a three-dimensional kernel the quadratic has rank
three, and each conic has a \(K\)-point. Lean checks the kernel
coordinates, the arithmetic tests, and the counts.

\begin{center}
\small
\begin{tabular}{@{}rll@{}}
\toprule
Matchings & Restricted quadratic & Matching locus over \(K\)\\
\midrule
36 & Two distinct rational roots & \(\operatorname{Spec}(K\times K)\)\\
36 & A double rational root & \(\operatorname{Spec}(K[\epsilon]/(\epsilon^2))\)\\
28 & An irreducible binary quadratic & \(\operatorname{Spec}(K_2)\)\\
5  & A nonsingular ternary quadratic & A smooth conic with a \(K\)-point\\
\bottomrule
\end{tabular}
\end{center}

Keeping the matching labels, the zero-dimensional part has coordinate
algebra
\[
 A_{\mathrm{fin}}\simeq K^{72}\times
       \bigl(K[\epsilon]/(\epsilon^2)\bigr)^{36}\times K_2^{28}.
\]
Its reduction is \(K^{108}\times K_2^{28}\), so it is finite \'etale over
\(K\). Each conic is isomorphic to \(\mathbb P^1_K\) by projection from
its rational point.

Over \(q=17^e\), the number of rational matching--pencil pairs is
\[
 108+56\,\mathbf 1_{2\mid e}+5(q+1)
\,.
\]
A pencil can occur with several matching labels.

\begin{proposition}[Uniform spectrum on the eight-site domain]
\label{prop:seventeen-full-spectrum}
Keep the same evaluation sites and set \(q=17^e\), \(e\ge1\). The
worst pair-MCA counts for the \([8,4]\) code over \(\mathbb F_q\) are
\begin{gather*}
 B_{\mathcal C}(0)=1,\qquad B_{\mathcal C}(1)=2,\qquad B_{\mathcal C}(2)=7,\\
 B_{\mathcal C}(r)=\min(q,56)\qquad(r\ge3).
\end{gather*}
\end{proposition}

\begin{proof}
The zero-radius count follows by testing a line through the zero
syndrome with nonzero direction. At count one each hit comes from a
single parity column, and any three columns are independent. A pencil
through two column directions realizes two hits with its point at
infinity elsewhere.

At count two, let \(W\) be a two-dimensional syndrome plane, and let
\(b\le2\) be the number of parity columns it contains. Partition the
other columns by their projective images in \(V/W\). Each class has size
at most \(3-b\), since it and the \(b\) columns lie in a three-dimensional
space and any four columns are independent. A pair of nonbase columns
contributes one hit exactly when they lie in the same class. Two such
pairs give distinct hits. Disjoint pairs span disjoint vector subspaces,
and pairs sharing a site meet only in the direction of that column, which
lies outside \(W\). None of these hits equals a base-column direction,
by independence of any three columns. Pairs containing a base column
give at most a direction already counted, and pairs of two base columns
span \(W\). Hence, with class sizes \(g_j\),
\[
 N_2(W)=b+\sum_j\binom{g_j}{2}\le7.
\]
For \(b=0\) the bound is attained by class sizes \(3,3,2\), and for
\(b=1,2\) the bounds are \(4,2\). As an explicit plane, take the common
kernel of the two rows represented by \(A(X)=X^2+10X+2\) and
\(B(X)=X^3-X+6\), using coefficient order from constant through cubic.
On the eight sites in the listed support order the ratios \(A(a)/B(a)\)
are \(5,13,13,\infty,13,\infty,5,5\), with no base columns. A point at
infinity outside the seven hits gives seven affine challenges. This
plane and its three classes persist in all extensions. Dependent
syndrome columns contribute at most one at both counts.

At count three, write \(R_E\) for the coefficient row of the monic
cubic with roots at the sites in a triple \(E\). This annihilates the
three parity columns, so the support span is \(\ker R_E\). Whenever
\(R_Et\ne0\), the corresponding bad parameter is \(-R_Es/R_Et\).
For \(e=1\), the columns \(s=(14,10,9,10)\), \(t=(12,12,2,2)\)
give every element of \(\mathbb F_{17}\) after dropping triples with
zero denominator. For \(e=2\), take
\begin{align*}
 s&=(15+16x,\ 8+x,\ 15x,\ 2),\\
 t&=(13+11x,\ 14+14x,\ 13+16x,\ 8+2x),\qquad x^2=11.
\end{align*}
All 56 denominators are nonzero, and the resulting ratios are pairwise
distinct, as checked by cross-multiplication in \(\mathbb F_{17}[x]/(x^2-11)\).
For \(e\ge3\), choose \(t=(0,0,0,1)\), so every denominator is one.
Avoid the \(\binom{56}{2}\) proper hyperplanes \((R_E-R_{E'})s=0\).
They cannot cover the syndrome space because \(q>\binom{56}{2}\).
These constructions attain \(\min(q,56)\), which is the upper bound
at every radius from \cref{thm:pair-mca}. Monotonicity gives the same
value for \(r\ge3\).
\end{proof}

\paragraph{The fractional cut across the tower.}
For \(e\le31\), the 128-bit budget is below one, so there is no safe
count for either problem. At \(e=32\) the budget allows six, so
\cref{prop:seventeen-full-spectrum} gives MCA counts \((1,2)\).
For lists the half-distance bound gives safety at count two. At count
three the syndrome \((1,0,0,1)\) belongs to every triple span whose sites
have product one. Multiplying all three sites by the same subgroup
element multiplies the product by its cube, and cubing permutes the
order-eight subgroup. Thus each product class has \(\binom{8}{3}/8=7\)
triples. This syndrome lies in no pair span, since the constant coefficient of its annihilating
quadratic is nonzero. Each triple therefore gives a distinct error
support, and embedding in one component gives the list counts \((2,3)\)
at every positive width. For \(e\ge33\) the budget is at least 118,
so MCA is safe through the full grid and the list counts are \((4,5)\)
by the capacity and next-shell theorems.

\paragraph{Motivic point counting.}
Let \(\mathcal I_{r,t}\) mark \(t\) ordered, distinct contributed points
on each \(\mathbb P(W)\), regardless of how many supports give a point.
Its fiber over \(W\) has \(\prod_{j=0}^{t-1}(N_r(W)-j)\) rational points.
Thus \(\mathcal I_{r,T+1}\) has point count zero exactly when the budget
\(1\le T<q\) holds. Point counting defines a motivic measure on the
Grothendieck ring of varieties \cite{ellenbergLarsen2010counting}.
In the four-point example the base is a point and all four directions
are rational, so \([\mathcal I_{1,t}]\) is \(\prod_{j=0}^{t-1}(4-j)\)
times the point class. For larger codes, stratifying the intersection
patterns could give these counts across field extensions.
